\documentclass{article}

% NeurIPS 2025 template — download neurips_2025.sty from:
% https://neurips.cc/Conferences/2025/PaperInformation/StyleFiles
% and place it in the same directory as this file.
\usepackage[final,nonatbib]{neurips_2024}

\usepackage{amsmath}
\usepackage{amssymb}
\usepackage{graphicx}
\usepackage{booktabs}
\usepackage{microtype}
\usepackage[hidelinks]{hyperref}
\usepackage{xspace}

\newcommand{\daf}{\textsc{daf}\xspace}

\title{Pre-trained Surrogates Suffice: Active Learning for Crystal Stability
  on the WBM Benchmark}

\author{Anonymous}

\begin{document}

\maketitle

\begin{abstract}
Iterative active learning (AL) can dramatically reduce the number of expensive
density functional theory (DFT) evaluations needed to find stable materials, but
most prior work evaluates on small synthetic datasets.
We apply AL to the full WBM dataset (256{,}963 crystal structures with DFT
stability labels) from the Matbench Discovery benchmark---the largest publicly
available test bed for crystal stability prediction---and report two
counter-intuitive findings.
First, \emph{fine-tuning the surrogate hurts}: a frozen pre-trained CHGNet
surrogate (trained on 700K Materials Project structures) outperforms one
fine-tuned on the small, biased labeled set, which suffers catastrophic
forgetting.
Second, \emph{greedy exploitation matches uncertainty-aware UCB}: with a
0.9\% labeling budget, both strategies achieve a
Discovery Acceleration Factor (\daf) of $1.13\pm 0.02$ versus
$1.00\pm 0.03$ for random---statistically indistinguishable from each other.
We argue that a well-calibrated foundation model renders the uncertainty signal
redundant at this budget scale.
All code and instructions to reproduce are available at
\url{https://github.com/0xSoftBoi/gnome-materials}.
\end{abstract}

\section{Introduction}

Computational materials discovery is bottlenecked by DFT: evaluating a single
crystal structure costs minutes to hours of CPU time, making exhaustive screening
of large candidate pools infeasible.
Active learning addresses this by querying only the most informative or promising
candidates at each iteration, reducing the labeling budget by orders of magnitude.

Existing AL studies in materials science operate on datasets of at most tens of
thousands of structures and typically train surrogate models from
scratch~\cite{merchant2023gn,xie2018crystal}.
The recently released WBM dataset~\cite{riebesell2024matbench}---256{,}963
crystal structures with DFT hull-distance labels---opens the door to large-scale
evaluation, and the Matbench Discovery paper explicitly notes iterative AL as an
open direction.

We fill this gap with three contributions:
\begin{enumerate}
  \item \textbf{WBM as an iterative AL benchmark.}
    We define a standard AL protocol on WBM (DAF metric, fixed budget) and
    release code to reproduce it.
  \item \textbf{No-fine-tuning finding.}
    Updating CHGNet~\cite{deng2023chgnet} on the small labeled set degrades
    performance.  Using the pre-trained weights frozen throughout is strictly
    better.
  \item \textbf{Greedy = UCB at 0.9\% budget.}
    MC Dropout uncertainty does not improve over pure exploitation when the
    surrogate is already well-calibrated for the candidate pool.
\end{enumerate}

\section{Method}

\paragraph{Surrogate.}
We use CHGNet~v0.3 (412{,}525 parameters)~\cite{deng2023chgnet}, a universal
neural network potential pre-trained on 700K Materials Project structures.
The graph-convolutional backbone is frozen; a dropout probability of $p{=}0.3$
is injected into the MLP readout head to enable stochastic inference.

\paragraph{Uncertainty estimation.}
We use MC Dropout~\cite{gal2016dropout}: the backbone runs in \texttt{eval()}
mode while the MLP head remains in \texttt{train()}, giving stochastic forward
passes.  With $T{=}10$ passes we obtain per-structure mean $\mu(x)$ and standard
deviation $\sigma(x)$.

\paragraph{Two-stage shortlist.}
Running $T$ forward passes on all 250K+ unlabeled structures per iteration is
prohibitively slow.
We instead use a two-stage approach:
(1)~a single deterministic forward pass on a random sample of
$\min(5k_\text{list},\,|U|)$ unlabeled candidates ranks them by $\mu$;
(2)~MC Dropout is applied only to the top-$k_\text{list}$ by $\mu$.
This concentrates uncertainty estimation on the most promising region of chemical
space and avoids the failure mode of UCB selecting highly uncertain but
energetically unfavorable structures.

\paragraph{Acquisition functions.}
\begin{itemize}
  \item \textbf{Random}: uniform random selection (baseline).
  \item \textbf{Greedy}: select the $k$ structures with lowest $\mu(x)$.
  \item \textbf{UCB}: select by $\mu(x) - \lambda\sigma(x)$, $\lambda{=}1.0$.
\end{itemize}
For minimization (lower formation energy = more stable), UCB exploration
corresponds to $-\lambda\sigma$ rather than the $+\beta\sigma$ convention in
maximization BO~\cite{auer2002using}.

\paragraph{No fine-tuning.}
At each iteration we set \texttt{epochs\_per\_iter}$=0$: the surrogate weights
are never updated.  We verified in preliminary experiments that fine-tuning on
the small, biased initial labeled set ($N \le 400$) catastrophically degrades
CHGNet's predictions---a consequence of the surrogate's prior (700K MP
structures) being far broader than the labeled subset.

\section{Experiments}

\paragraph{Dataset.}
The WBM dataset~\cite{riebesell2024matbench} contains 256{,}963 initial crystal
structures with DFT hull-distance labels
(\texttt{e\_above\_hull\_mp2020\_corrected\_ppd\_mp}).
A structure is \emph{stable} if its hull distance $\le 0$\,eV/atom; 42{,}825
structures qualify (16.7\% prevalence).

\paragraph{Protocol.}
We start with 200 randomly selected labeled structures and run 20 iterations of
100 labels each, for a total budget of 2{,}200 (0.9\% of the pool).
Experiments are repeated for 5 random seeds and we report mean $\pm$ std.

\paragraph{Metric.}
The Discovery Acceleration Factor (\daf) measures how much more efficiently a
strategy finds stable structures relative to random:
\[
  \daf = \frac{n_\text{stable} / B}{n_\text{stable,total} / N},
\]
where $B$ is the labeling budget, $N$ the pool size, and
$n_\text{stable,total}=42{,}825$.
A random strategy achieves $\daf\approx 1.0$; a perfect oracle achieves
$\daf = N / n_\text{stable,total} \approx 6.0$.

\section{Results}

\begin{table}[t]
  \centering
  \caption{WBM active learning results (mean $\pm$ std across 5 seeds).
           Budget: 2{,}200 labels (0.9\% of 256{,}963 structures).}
  \label{tab:wbm}
  \begin{tabular}{lccc}
    \toprule
    Strategy & Stable found & Precision & \daf \\
    \midrule
    Random          & $364.8 \pm 10.2$ & $16.6\% \pm 0.5\%$ & $0.995 \pm 0.028$ \\
    Greedy ($\mu$)  & $415.8 \pm 6.2$  & $18.9\% \pm 0.3\%$ & $1.134 \pm 0.017$ \\
    UCB ($\lambda{=}1$) & $414.4 \pm 9.5$ & $18.8\% \pm 0.4\%$ & $1.130 \pm 0.026$ \\
    \bottomrule
  \end{tabular}
\end{table}

Table~\ref{tab:wbm} and Figure~\ref{fig:results} summarise the main results.
Both Greedy and UCB find approximately 50 more stable structures than random
within the same 2{,}200-label budget---a 14\% improvement in stable-material
yield.
The key finding is that \textbf{Greedy and UCB are statistically
indistinguishable}: their \daf values overlap at one standard deviation.

The \daf curves in Figure~\ref{fig:results} (left) show both informed strategies
maintaining a stable advantage over random from iteration 2 onward, with no sign
of divergence between Greedy and UCB across the full campaign.

\begin{figure}[t]
  \centering
  \includegraphics[width=\linewidth]{results/wbm_multiseed.png}
  \caption{WBM active learning results across 5 seeds (mean $\pm 1\sigma$ shaded).
    \textbf{Left:} Discovery Acceleration Factor vs.\ labeling budget.
    Dashed line: random baseline (\daf$=1$); dotted line: perfect oracle
    (\daf$\approx 6.0$).
    \textbf{Right:} Cumulative stable structures found.
    Greedy and UCB are indistinguishable throughout; both outperform random
    from iteration 2.}
  \label{fig:results}
\end{figure}

\paragraph{Why Greedy = UCB?}
Two mechanisms explain this result.
First, at 16.7\% prevalence, the top-$\mu$ region of CHGNet's predictions is
already highly enriched for stable structures: the Stage~1 filter (lowest
predicted $\mu$) acts as a strong positive signal on its own, leaving little
room for $\sigma$ to improve precision further.
Second, CHGNet's backbone is frozen, so $\sigma$ estimates reflect only readout
head uncertainty---which is nearly uniform across the diverse WBM chemistry.
A fine-tuned backbone could produce more informative, structure-specific $\sigma$
values, but fine-tuning hurts $\mu$ accuracy (see below).

\paragraph{Synthetic scaling context.}
On a 50K-structure synthetic pool with a GCN surrogate trained from scratch, we
find UCB outperforms Greedy by 25 percentage points in top-10 recall (70\% vs.\
45\%).  The contrast with WBM highlights that the value of uncertainty
exploration depends critically on whether the surrogate is already well-calibrated
for the candidate distribution.  Pre-trained foundation models collapse the
exploitation--exploration gap.

\section{Discussion and Conclusion}

\paragraph{No fine-tuning.}
CHGNet's pre-training on 700K diverse MP structures provides a broad prior over
crystal chemistry.  Fine-tuning on $\le 400$ WBM structures selected by an early,
biased policy shifts the readout MLP toward a narrow region of chemical space,
degrading generalisation to the remaining 256K candidates.
This is consistent with the catastrophic forgetting literature~\cite{goodfellow2013empirical}
and suggests that \emph{foundation model surrogates should be used frozen} in
iterative AL unless the fine-tuning is regularised (e.g., elastic weight
consolidation).

\paragraph{Implications for AL design.}
Our results suggest a practical recipe for large-scale materials AL: (1)~use a
pre-trained universal potential frozen, (2)~apply a two-stage shortlist to make
MC Dropout tractable on 250K+ structures, and (3)~use Greedy acquisition---UCB
adds complexity without benefit when the surrogate is already calibrated.

\paragraph{Limitations and future work.}
Our budget (0.9\%) is small relative to the pool; at larger budgets the
exploitation--exploration trade-off may reassert itself.
We evaluate a single surrogate family (CHGNet); extending to M3GNet or TensorNet
via the MatGL library~\cite{matgl2024} would test generality.
The WBM benchmark is one-shot classification in its standard formulation;
our iterative protocol is a complementary framing, not a replacement.

\paragraph{Conclusion.}
We demonstrate that iterative active learning on the 256K WBM benchmark achieves
a consistent 14\% improvement in stable-material yield over random screening
at 0.9\% budget, using a frozen pre-trained CHGNet surrogate with MC Dropout.
The surprising finding---that greedy exploitation matches uncertainty-aware
UCB---points to a general principle: when the surrogate is a well-calibrated
foundation model, uncertainty estimation provides no additional gain at low
labeling budgets.

\begin{thebibliography}{99}

\bibitem{riebesell2024matbench}
Riebesell, J., et al.\ (2024).
Matbench Discovery: A framework to evaluate machine learning crystal stability
predictions.
\textit{Nature Machine Intelligence}.
\url{https://doi.org/10.1038/s42256-025-01055-1}

\bibitem{deng2023chgnet}
Deng, B., Zhong, P., Jun, K., Riebesell, J., Han, K., Bartel, C.\ J.,
Ceder, G.\ (2023).
CHGNet as a pretrained universal neural network potential for
charge-informed atomistic modelling.
\textit{Nature Machine Intelligence}, 5, 1031--1041.

\bibitem{merchant2023gn}
Merchant, A., Batzner, S., Schoenholz, S.\ S., et al.\ (2023).
Scaling deep learning for materials discovery.
\textit{Nature}, 624, 80--85.

\bibitem{gal2016dropout}
Gal, Y., Ghahramani, Z.\ (2016).
Dropout as a Bayesian approximation: Representing model uncertainty in deep
learning.
\textit{Proceedings of ICML}, 1050--1059.

\bibitem{auer2002using}
Auer, P., Cesa-Bianchi, N., Fischer, P.\ (2002).
Finite-time analysis of the multiarmed bandit problem.
\textit{Machine Learning}, 47(2), 235--256.

\bibitem{xie2018crystal}
Xie, T., Grossman, J.\ C.\ (2018).
Crystal graph convolutional neural networks for an accurate and interpretable
prediction of material properties.
\textit{Physical Review Letters}, 120(14), 145301.

\bibitem{goodfellow2013empirical}
Goodfellow, I.\ J., Mirza, M., Xiao, D., Courville, A., Bengio, Y.\ (2013).
An empirical investigation of catastrophic forgetting in gradient-based neural
networks.
\textit{arXiv:1312.6211}.

\bibitem{matgl2024}
Materialyze.AI and Intel Labs.\ (2024).
MatGL: Materials Graph Library.
\url{https://github.com/materialyzeai/matgl}

\end{thebibliography}

\end{document}
